% Nouveau EGA V (cote 25): archive cote 25, site page 39 (PDF page 40).
% typescript. Draft compared: none.
% AI-transcribed by Claude Opus 5.5 (claude-opus-5-5, adaptive thinking, effort high), Message Batches API, batch ega25b (msgbatch_01A7n5YuXubUPNK1zsepi8RG), 2026-11-13; original AI transcription, subsequently checked against both supplied scan parts; see PAGES_CHECKED.tsv and CHECK_DECISIONS.tsv.
\scanpage{25/39}
\mspage{\hand{\unsure{35|33}}}

une sous-variété linéaire passant par a son espace tangent en a, consi-%
déré comme sous-espace  de l'espace tangent à P en a. Ce genre
\placedmargin{\hand{||}}de \unsure{sorites|sorttes}, ensemble avec divers sorites sur les sous-variétés linéaires
et les grassmaniennes, devraient être donnés dans un ou deux n$^{\text{o}}$s prélimi
naires, bien entendu en énonçant tout systématiquement sur une base
quelconque.

En fait, il y a mieux, savoir que le \uline{sous-préschéma} $\underline{Y}^{\mathrm{sing}}$
des zéros singuliers de $\hand{\phi}$ relativement à $\overset{\hand{\vee}}{P}$ , défini dans par.16, n'est
autre que $P(\hand{\unsure{\underline{\nu}|\check{\underline{\nu}}}}_{X/P})$ , et (comme ce dernier est lisse sur S de dimension
relative d \unsure{$+$} (r-d-1) = r - 1 (r étant la dim relative de P sur S)
on est sous les conditions fav\unsure{o|n}rables étudiées au par.16 .  Pour le vérifier,
notons que par définition, $\underline{Y}^{\mathrm{sing}}$ n'est autre que le sous-préschéma de $\underline{Y}$
des zéros de la section $\psi =$ $d\hand{\phi}|\underline{Y}$ de $\hand{\underline{\Omega}}^1_{X_{\overset{\hand{\vee}}{P}}/\overset{\hand{\vee}}{P}}(1,1)\otimes \underline{O}_Y = \hand{\underline{\Omega}}^1_{X/S}\otimes\underline{O}_Y(1,1)$\hand{.}
Nous allons donner une autre interprétation de cette section, dont
la conclusion résultera aussitôt. Pour ceci, considérons le diagramme
 suites exactes
suivant sur $X_{\overset{\hand{\vee}}{P}}$ , ou plus généralement sur un préschéma Z quelconque au dessus de
$X_{\overset{\hand{\vee}}{P}}$ :

\mspage{\hand{39}}
